% !TEX program = xelatex
\documentclass{misisarticle}

\title{Оптимизация кредитных решений на основе оценки вероятности дефолта методами машинного обучения с использованием критерия чистой приведённой стоимости}
\misisauthor{Григорян Виктория}{НИТУ МИСИС, прикладная математика}{}{}

\abstract{Исследуется влияние ошибок оценки вероятности дефолта на выбор кредитных заявок при ограниченном бюджете. Цель работы~--- сформулировать задачу максимизации ожидаемой чистой приведённой стоимости портфеля и описать методику её решения с использованием логистической регрессии и градиентного бустинга. В однопериодной модели получена точная зависимость ошибки оценки стоимости кредита от ошибки прогноза вероятности дефолта. Выбор заявок сформулирован как задача бинарного линейного программирования. Предложен протокол сравнения исходных и откалиброванных моделей на отложенной выборке по метрикам ранжирования, вероятностных прогнозов и стоимости выбранных портфелей в общих сценариях денежных потоков.}

\keywords{кредитный скоринг, вероятность дефолта, машинное обучение, логистическая регрессия, градиентный бустинг, математическая оптимизация, чистая приведённая стоимость}

\begin{document}
\maketitle

\section{Введение}

Кредитный скоринг используется для систематической оценки риска невозврата по заявкам на потребительские кредиты. В исторических данных исход кредита связывается с характеристиками заёмщика и условиями договора, известными к моменту принятия решения. Оценка, построенная по этим данным, может использоваться для ранжирования заявок, оценки вероятности дефолта (PD, \textit{probability of default}) и формирования правил одобрения. Обзор статистических методов скоринга дан Хэндом и Хенли~\cite{hand}, а сравнительное исследование классических моделей и ансамблей представлено Лессманном и соавторами~\cite{lessmann}.

Способность модели выделять более рискованные заявки измеряется, в частности, площадью под ROC-кривой. Это свойство зависит от порядка оценок. Строго возрастающее преобразование прогнозов сохраняет ROC-AUC, хотя численные вероятности после преобразования могут существенно отличаться. Если среди заявок с прогнозом $0{,}2$ доля дефолтов систематически отличается от $0{,}2$, то расчёт ожидаемых поступлений по этим кредитам получает систематическое смещение. Калибровка направлена на согласование прогнозируемых вероятностей с наблюдаемыми частотами исходов~\cite{niculescu}.

Экономические последствия ошибок зависят от условий конкретного договора. Для двух заявок с одинаковой PD могут различаться сумма кредита, возврат при исполнении обязательств и поступления при дефолте. При ограниченном бюджете выдача одного кредита также сокращает объём ресурсов, доступных для других заявок. Поэтому результаты ранжирования требуется связать с денежными потоками и задачей выбора портфеля. Такой переход согласуется с исследованиями, в которых скоринговые модели оцениваются по последствиям кредитных решений~\cite{verbraken,lessmann}.

Исследовательский вопрос состоит в том, как качество оценки PD влияет на выбор заявок по критерию максимизации ожидаемой чистой приведённой стоимости (NPV). Цель работы~--- разработать математическую постановку этой задачи и описать методы её решения на основе логистической регрессии и градиентного бустинга, предусмотрев независимую оценку выбранных портфелей.

\section{Анализ существующих работ}

\subsection{Кредитный скоринг и статистические модели}

Хэнд и Хенли~\cite{hand} систематизируют методы классификации для потребительского скоринга и рассматривают трудности, связанные с формированием обучающей выборки. В частности, при наличии исходов только по ранее одобренным заявкам связь признаков с риском оценивается на выборке, сформированной прежней кредитной политикой. Перенос оценки на всех новых заявителей требует дополнительных оснований. Это ограничение существенно и при экономической оценке модели, поскольку оптимизация может выбирать заявки в областях пространства признаков, слабо представленных в наблюдениях.

В логистической регрессии коэффициент при признаке задаёт изменение логарифма шансов дефолта при единичном увеличении признака и фиксированных остальных переменных. Такая параметризация облегчает проверку направления связей и делает модель естественной точкой отсчёта. Линейность логарифма шансов ограничивает представление пороговых эффектов и взаимодействий. Для их учёта требуется явно включать соответствующие преобразования и произведения признаков.

Лессманн и соавторы~\cite{lessmann} сравнивают классификаторы на нескольких наборах кредитных данных, применяя статистические и стоимостные критерии. Для настоящей работы важен сам принцип сопоставления моделей по разным целям. Выводы сравнительного исследования зависят от данных, метрик и правил принятия решений. В частности, оценка полезности при заданной структуре стоимости ошибок не определяет состав оптимального портфеля при неоднородных суммах кредитов и общем бюджете.

\subsection{Градиентный бустинг в задачах классификации}

Фридман~\cite{friedman} описывает построение аддитивной модели как последовательное приближение направления спуска в пространстве функций. В качестве базовой модели можно использовать дерево, которое разбивает пространство признаков на области с разными значениями прогноза. Разбиения на нескольких уровнях позволяют выразить взаимодействия. Например, эффект размера кредита может зависеть от уровня дохода, если порог по размеру кредита используется только в одной ветви разбиения по доходу.

Чен и Гестрин~\cite{xgboost} развивают этот подход в XGBoost, используя регуляризацию сложности деревьев и приближение целевой функции второго порядка при их построении. Эти средства помогают управлять гибкостью ансамбля и вычислительными затратами. Для запланированного сравнения достаточно стандартного градиентного бустинга над неглубокими деревьями. Работа о XGBoost служит примером развития метода, а её результаты по масштабируемости не характеризуют калибровку PD на выбранных кредитных данных.

\subsection{Калибровка вероятностей}

Никулеску-Мизил и Каруана~\cite{niculescu} исследуют качество вероятностей, получаемых от разных алгоритмов обучения, и сопоставляют сигмоидальную и изотоническую калибровку. Сигмоидальная калибровка задаёт параметрическое преобразование исходной оценки, тогда как изотоническая оценивает произвольную неубывающую зависимость. Дополнительная гибкость изотонической регрессии увеличивает требования к объёму калибровочных данных. Это особенно существенно при малом числе наблюдений в крайних диапазонах риска.

Идеальная калибровка прогнозов означает $\mathbb{E}[Y\mid\hat p]=\hat p$. Данное свойство относится к группам с одинаковым прогнозом. Оно допускает ошибки для отдельных заёмщиков и для подгрупп, например при разных размерах кредита. В задаче оптимизации такие ошибки могут иметь неодинаковые последствия. При одинаковой ошибке PD изменение NPV больше по кредиту с большим разрывом между поступлениями при исполнении договора и при дефолте. Поэтому калибровку требуется анализировать вместе с экономическими характеристиками заявок.

\subsection{Экономически ориентированные кредитные решения}

Вербракен и соавторы~\cite{verbraken} предлагают критерий ожидаемой максимальной прибыли для оценки скоринга. Критерий учитывает потери от выдачи кредита заёмщику высокого риска и упущенный доход при отказе заёмщику, который исполнил бы обязательства. Авторы используют стоимостный критерий для выбора порога скоринговой модели. Здесь исследуется другое правило принятия решений. Для каждой заявки задаются собственные денежные потоки, а одобрения связаны общим бюджетом. Единый порог по PD в общем случае не решает полученную задачу.

Таким образом, сравнение классификаторов по статистическим критериям оставляет открытым вопрос о влиянии прогнозов PD на результат последующей оптимизации. Для его изучения требуется применить одинаковые денежные сценарии и бюджет ко всем моделям, а затем оценить выбранные портфели на независимых данных.

\section{Постановка задачи}

\subsection{Оценка вероятности дефолта}

Пусть $\mathcal D=\{(x_i,y_i)\}_{i=1}^{N}$~--- обучающая выборка, $x_i\in\mathbb R^d$~--- вектор признаков после кодирования, $y_i\in\{0,1\}$~--- наблюдаемое значение индикатора дефолта $Y_i$. Значение $1$ соответствует дефолту в заданном горизонте. Обозначим истинную условную вероятность через $p_i=\mathbb P(Y_i=1\mid X_i=x_i)$, а прогноз по модели через $\hat p_i$. Цель обучения состоит в приближении $p_i$ с использованием признаков, доступных перед выдачей кредита.

\subsection{Логистическая регрессия}

Пусть $\sigma(t)=(1+\exp(-t))^{-1}$. Модель задаётся формулой
\begin{equation}
  \hat p_i=\sigma(\beta_0+\beta^\mathsf{T}x_i),
  \label{eq:logit}
\end{equation}
где $\beta_0\in\mathbb R$~--- свободный член, $\beta\in\mathbb R^d$~--- вектор коэффициентов. При $\ell(y,q)=-y\log q-(1-y)\log(1-q)$ параметры определяются из задачи
\begin{equation}
  \min_{\beta_0,\beta}\quad
  \frac{1}{N}\sum_{i=1}^N\ell(y_i,\hat p_i)
       +\lambda\|\beta\|_2^2 .
  \label{eq:training}
\end{equation}
Первое слагаемое представляет среднее отрицательное логарифмическое правдоподобие для бернуллиевских исходов. Оно сильно штрафует прогнозы, назначающие малую вероятность наблюдаемому исходу. Второе слагаемое с $\lambda\geq 0$ ограничивает величину коэффициентов. При увеличении $\lambda$ возрастают издержки больших коэффициентов, что позволяет контролировать переобучение. Свободный член обычно оставляют без штрафа, чтобы модель могла учитывать базовую частоту дефолта.

Оптимизация в~\eqref{eq:training} непрерывна. Если $t=\beta_0+\beta^\mathsf{T}x_i$, то $\ell(y_i,\sigma(t))=\log(1+\exp t)-y_i t$. Вторая производная по $t$ равна $\sigma(t)(1-\sigma(t))\geq 0$. Композиция с аффинным аргументом сохраняет выпуклость, а квадратичный штраф при $\lambda\geq0$ также выпуклый.

\subsection{Градиентный бустинг}

В бустинге~\cite{friedman} логарифм шансов представлен аддитивной моделью
\begin{equation}
  \begin{gathered}
    F_M(x)=F_0(x)+\sum_{m=1}^{M}\eta f_m(x),\\
    \hat p(x)=\sigma(F_M(x)),
  \end{gathered}
  \label{eq:boost}
\end{equation}
где $F_0$~--- начальная модель, $M$~--- число итераций, $\eta>0$~--- шаг обучения, $f_m$~--- решающее дерево. Для логистических потерь отрицательный градиент по текущей оценке $F_{m-1}(x_i)$ равен $y_i-\sigma(F_{m-1}(x_i))$. На очередном шаге дерево приближает этот вектор на обучающих наблюдениях, после чего его вклад добавляется с множителем $\eta$. Ограничение глубины деревьев и шага регулирует сложность ансамбля. Поиск разбиений выполняется жадно, а совместная подгонка структур всех деревьев включает дискретный выбор и не образует выпуклой задачи.

\subsection{Ожидаемая чистая приведённая стоимость}

Пусть $L_i>0$~--- объём выдачи в начале периода, $R_i$~--- поступление в конце периода при исполнении договора, $D_i$~--- поступление при дефолте, а $r>-1$~--- заданная ставка дисконтирования за период. Для произвольной условной вероятности дефолта $q$ ожидаемая стоимость равна
\begin{equation}
  V_i(q)=-L_i+\frac{(1-q)R_i+qD_i}{1+r}.
  \label{eq:npv}
\end{equation}
В оптимизации используется оценка $\widehat{\mathrm{NPV}}_i=V_i(\hat p_i)$. Сумма выдачи учитывается в момент принятия решения, а условные поступления взвешиваются по вероятностям взаимоисключающих исходов и дисконтируются.

Модель рассматривает один период и бинарный исход. Условные денежные потоки считаются известными, ставка задаётся заранее. Транзакционные и операционные издержки отдельно не моделируются. Упрощение позволяет выделить влияние ошибки PD при фиксированных денежных потоках, оставляя за рамками работы график платежей, момент дефолта внутри периода и динамику взыскания.

\subsection{Чувствительность NPV к ошибке PD}

Из~\eqref{eq:npv} следуют производная и точная формула для изменения стоимости
\begin{align}
  \frac{\partial V_i(q)}{\partial q}
     &=\frac{D_i-R_i}{1+r}, \label{eq:derivative}\\
  V_i(\hat p_i)-V_i(p_i)
     &=\frac{D_i-R_i}{1+r}(\hat p_i-p_i).
     \label{eq:error}
\end{align}
При $R_i>D_i$ завышение PD уменьшает оценочную NPV относительно $V_i(p_i)$, а занижение увеличивает её. Поскольку зависимость аффинна, формула~\eqref{eq:error} справедлива и для конечных отклонений $\Delta p_i=\hat p_i-p_i$. Модуль ошибки пропорционален дисконтированной разности поступлений по двум исходам.

При отсутствии бюджетной конкуренции заявка с положительной NPV выгодна в рамках модели. Для $R_i>D_i$ условие $V_i(q)>0$ эквивалентно $q<q_i^\star$, где
\begin{equation}
  q_i^\star=\frac{R_i-(1+r)L_i}{R_i-D_i}.
  \label{eq:threshold}
\end{equation}
Порог зависит от условий кредита и может лежать вне диапазона $[0,1]$. При $0<q_i^\star<1$ ошибка прогноза способна изменить знак NPV, если прогноз и истинная вероятность лежат по разные стороны порога. При активном бюджетном ограничении изменение состава портфеля возможно и при сохранении знаков NPV отдельных заявок.

\subsection{Оптимальный выбор кредитного портфеля}

Для пакета из $n$ заявок введём $z_i\in\{0,1\}$, где $z_i=1$ означает одобрение кредита в полном объёме. При бюджете $B\geq0$ задача имеет вид
\begin{equation}
  \begin{aligned}
    \max_{z}\quad &\sum_{i=1}^{n}z_i V_i(\hat p_i),\\
    \text{при}\quad &\sum_{i=1}^{n}z_i L_i\le B,\\
    &z_i\in\{0,1\},\qquad i=1,\ldots,n .
  \end{aligned}
  \label{eq:knapsack}
\end{equation}
Целевая функция суммирует ожидаемые NPV одобренных кредитов. Ограничение учитывает общий объём средств, размещаемых при выдаче. Бинарные переменные соответствуют неделимым решениям об одобрении. После фиксации прогнозов и денежных параметров получается задача бинарного линейного программирования со структурой задачи о рюкзаке. Оценивание коэффициентов логистической регрессии относится к непрерывной оптимизации, а выбор портфеля~--- к дискретной.

Аддитивность математического ожидания стоимости портфеля следует из линейности ожидания и не требует независимости дефолтов. Зависимость между исходами имеет значение для оценки вариации потерь и вероятности крупных убытков. Задача~\eqref{eq:knapsack} таких совместных ограничений на портфельный риск не содержит.

Для оценки влияния ошибок на выбор можно сравнить портфель $\hat z$, оптимальный при $\hat p$, с портфелем $z^\star$, оптимальным при $p$. Пусть $S(z,q)=\sum_i z_i V_i(q_i)$. Тогда
\begin{equation}
  \begin{aligned}
    0 &\le S(z^\star,p)-S(\hat z,p)\\
      &\le \sum_{i=1}^n |z_i^\star-\hat z_i|
        \frac{|R_i-D_i|}{1+r}|\hat p_i-p_i|.
  \end{aligned}
  \label{eq:regret}
\end{equation}
Для верхней оценки к разности портфельных NPV прибавляются и вычитаются стоимости обоих портфелей, рассчитанные при $\hat p$. Разность $S(z^\star,\hat p)-S(\hat z,\hat p)$ неположительна по оптимальности $\hat z$. Оценивание оставшейся суммы по модулю с помощью~\eqref{eq:error} даёт~\eqref{eq:regret}. В правой части участвуют только заявки, по которым решения расходятся. Поскольку истинные PD неизвестны, эта граница объясняет механизм потерь и используется для интерпретации анализа чувствительности.

\section{Описание методов}

\subsection{Исходные данные}

Для вычислительного исследования предлагается Statlog (German Credit Data)~\cite{uci}. Исходная версия \texttt{german.data} содержит 1000 наблюдений и 20 признаков, в том числе кредитную сумму и продолжительность договора. Документация обозначает классы как \textit{good} и \textit{bad}, с кодами $1$ и $2$ соответственно. Перед обучением следует сверить коды и трактовку признаков с описанием данных, в том числе с исправленной документацией South German Credit~\cite{south}. Для учебной постановки \textit{bad} предлагается отобразить в $y=1$, а \textit{good}~--- в $y=0$. Метка \textit{bad} используется как приближение дефолта, поскольку исходная документация не устанавливает согласованного горизонта наблюдения и операционного определения дефолта для банковской оценки PD.

Набор не содержит полных потоков платежей, условий взыскания и данных, необходимых для проверки фактической NPV. Предлагается принять задокументированную кредитную сумму за $L_i$, а остальные параметры варьировать по общим сценариям. Например, можно задать $R_i=(1+a)L_i$, $D_i=\rho L_i$, где $a$~--- условная надбавка к сумме выдачи, а $\rho\in[0,1]$~--- доля возврата в случае дефолта. Если надбавка и доля возврата одинаковы для всех заявок, сценарий отражает неоднородность сумм при одинаковой относительной доходности. При задании нескольких условных групп условий можно изучить эффект неоднородности доходности и возврата. Правила образования групп фиксируются без использования тестовых исходов.

\subsection{Предобработка и валидация}

Исходную выборку предполагается однократно разделить на обучающую и тестовую части со стратификацией по метке. Одно и то же разбиение применяется ко всем моделям. Категориальные признаки кодируются индикаторами категорий, а числовые признаки для логистической регрессии стандартизируются, чтобы величина штрафа не зависела от выбора единиц измерения. Для деревьев стандартизация числовых признаков не требуется. Неизвестные при обучении категории обрабатываются заранее определённым правилом кодировщика. Словари категорий и параметры масштабирования оцениваются внутри обучающей части каждого разбиения.

Из состава признаков исключаются данные, недоступные в момент заявки. Гиперпараметры подбираются по стратифицированной перекрёстной проверке на обучающей части. Чтобы прогнозы для калибратора оставались вневыборочными даже относительно подбора гиперпараметров, при построении этих прогнозов настройку следует проводить внутри соответствующего обучающего фолда. Тестовые метки используются после фиксации моделей, способов калибровки и правил построения денежных сценариев.

\subsection{Обучение и калибровка}

Для логистической регрессии подбирается коэффициент регуляризации, для градиентного бустинга~--- глубина базовых деревьев, шаг обучения и число деревьев. В качестве общего критерия подбора можно использовать среднюю логарифмическую потерю на валидационных фолдах. Достаточно реализаций логистической регрессии и градиентного бустинга из \textit{scikit-learn}. Для воспроизводимости фиксируются разбиения, диапазоны гиперпараметров и начальные состояния генераторов случайных чисел.

Сигмоидальная калибровка строится как $g(s)=\sigma(as+b)$, где $s$~--- исходная оценка логарифма шансов, а параметры $a,b$ подбираются по вневыборочным прогнозам. При $a>0$ преобразование сохраняет ранжирование. Изотоническая регрессия ищет неубывающую функцию $g$, минимизирующую сумму квадратов отклонений от наблюдаемых меток. Она может присваивать одинаковые значения нескольким исходным оценкам, меняя обработку совпадений в ROC-AUC. На малой выборке изотоническая зависимость может переобучиться. Обе процедуры обучаются внутри обучающей части, исходные и откалиброванные модели оцениваются на одном тестовом наборе.

\subsection{Качество вероятностных прогнозов}

ROC-AUC интерпретируется как вероятность того, что случайный объект положительного класса получит более высокий балл, чем случайный объект отрицательного класса, с половинным весом совпадений. Brier Score и Log Loss на тестовой выборке $\mathcal T$ размера $n$ вычисляются как
\begin{align}
  \mathrm{BS} &= \frac{1}{n}\sum_{i\in\mathcal T}(y_i-\hat p_i)^2,\\
  \mathrm{LL} &= \frac{1}{n}\sum_{i\in\mathcal T}\ell(y_i,\hat p_i).
\end{align}
Обе метрики учитывают величину прогнозируемой вероятности. Логарифмические потери дополнительно чувствительны к уверенным ошибкам. При численном вычислении Log Loss для прогнозов $0$ и $1$ используется одна и та же заранее заданная отсечка для всех моделей. Диаграммы надёжности дополняют метрики сравнением средней PD с частотой положительной метки по интервалам прогнозов. Для интерпретации крайних интервалов следует показывать число наблюдений в них. Связь с оптимизацией даёт формула~\eqref{eq:error}, в которой величина PD непосредственно определяет оценочный денежный эффект.

\subsection{Решение задачи выбора портфеля}

Для каждого сценария денежных потоков и каждого набора прогнозов вычисляются коэффициенты $V_i(\hat p_i)$ и формируется задача~\eqref{eq:knapsack}. Сопоставимые бюджеты можно задать как фиксированные доли суммарного объёма заявок тестового пакета. Для минимизирующего интерфейса \texttt{scipy.optimize.milp} передаются коэффициенты $-V_i(\hat p_i)$, границы $0\le z_i\le 1$, условие целочисленности и одна бюджетная строка с коэффициентами $L_i$.

Метод ветвей и границ использует непрерывные релаксации для получения оценок оптимума и ветвится по переменным с дробными значениями. Допустимость выбранного портфеля проверяется по бинарности и фактической сумме выдач. Качество решения контролируется по статусу решателя и разрыву между значением найденного портфеля и границей. При сравнении портфелей фиксируются одинаковые параметры решателя, порядок заявок и бюджет. Это важно, поскольку несколько оптимальных портфелей могут иметь одинаковую NPV.

\subsection{План сравнения моделей и решений}

В каждом сценарии для исходных и откалиброванных моделей сравниваются ROC-AUC, Brier Score и Log Loss, число одобренных заявок, использованный бюджет, пересечение выбранных множеств и ожидаемая NPV по прогнозам. Сумма NPV, вычисленная из собственных PD модели, характеризует её прогноз для выбранных заявок. Эта сумма может вырасти из-за систематического занижения риска, поэтому её величина сама по себе не доказывает экономического превосходства модели.

Независимую оценку предлагается построить по отложенным меткам при неизменных сценарных денежных потоках. Пусть $C_i(0)=R_i$, $C_i(1)=D_i$. Сценарная стоимость выбранного портфеля при наблюдаемых исходах равна
\begin{equation}
  W(z,y)=\sum_{i\in\mathcal T}z_i
      \left(-L_i+\frac{C_i(y_i)}{1+r}\right).
  \label{eq:holdout}
\end{equation}
Значение $W$ отражает последствия решений в принятом учебном сценарии. Оно зависит от трактовки меток и предположения, что зарегистрированный исход применим к выдаче по заданным условиям. Для всех моделей используется один тестовый пакет, что позволяет сравнить разности $W$ при одних и тех же метках. При оценке неопределённости разностей можно повторно формировать парные тестовые пакеты из одних и тех же наблюдений для всех моделей и заново решать задачу при той же доле бюджета. Такой анализ требует предположения о репрезентативности наблюдений и не восстанавливает фактическую прибыль кредитной организации.

Для анализа чувствительности предполагается согласованно сдвигать прогнозы PD, ограничивая их интервалом $[0,1]$, и повторно выбирать портфель. Для каждого сдвига фиксируются число изменённых решений, бюджет и разность сценарных стоимостей~\eqref{eq:holdout}. Сопоставление этих показателей при разных допущениях о возврате и бюджете позволит установить, насколько выводы зависят от калибровки и от выбранных параметров денежных потоков.

\begin{thebibliography}{9}

\bibitem{hand}
Hand D. J., Henley W. E. Statistical classification methods in consumer credit scoring: a review. \textit{Journal of the Royal Statistical Society: Series A}, vol.~160, no.~3, pp.~523--541, 1997.
\url{https://doi.org/10.1111/j.1467-985X.1997.00078.x}.

\bibitem{lessmann}
Lessmann S., Baesens B., Seow H.-V., Thomas L. C. Benchmarking state-of-the-art classification algorithms for credit scoring: an update of research. \textit{European Journal of Operational Research}, vol.~247, no.~1, pp.~124--136, 2015.
\url{https://doi.org/10.1016/j.ejor.2015.05.030}.

\bibitem{niculescu}
Niculescu-Mizil A., Caruana R. Predicting good probabilities with supervised learning. \textit{Proceedings of the 22nd International Conference on Machine Learning}, pp.~625--632, 2005.
\url{https://doi.org/10.1145/1102351.1102430}.

\bibitem{verbraken}
Verbraken T., Bravo C., Weber R., Baesens B. Development and application of consumer credit scoring models using profit-based classification measures. \textit{European Journal of Operational Research}, vol.~238, no.~2, pp.~505--513, 2014.
\url{https://doi.org/10.1016/j.ejor.2014.04.001}.

\bibitem{friedman}
Friedman J. H. Greedy function approximation: a gradient boosting machine. \textit{The Annals of Statistics}, vol.~29, no.~5, pp.~1189--1232, 2001.
\url{https://doi.org/10.1214/aos/1013203451}.

\bibitem{xgboost}
Chen T., Guestrin C. XGBoost: a scalable tree boosting system. \textit{Proceedings of the 22nd ACM SIGKDD International Conference on Knowledge Discovery and Data Mining}, 2016.
\url{https://doi.org/10.1145/2939672.2939785}.

\bibitem{uci}
Hofmann H. Statlog (German Credit Data). \textit{UCI Machine Learning Repository}, 1994.
\url{https://archive.ics.uci.edu/dataset/144}.

\bibitem{south}
South German Credit. \textit{UCI Machine Learning Repository}.
\url{https://archive.ics.uci.edu/dataset/573}.

\end{thebibliography}
\end{document}
